\documentclass[10pt,a4paper]{amsart}
\usepackage[margin=19mm]{geometry}
\usepackage{amsmath,amssymb,mathtools}
\usepackage[scr=rsfs,cal=euler]{mathalfa}
\usepackage{xfrac}
\usepackage[unicode,hidelinks,bookmarks=false]{hyperref}
\newcommand{\ie}{i.e.\ }
\newcommand{\cf}{cf.\ }
\newcommand{\resp}{respectively\ }
\newcommand{\Subsection}{\subsection}
\providecommand{\nobreakdash}{\nobreak\hbox{-}}
\setlength{\emergencystretch}{2em}
\multlinegap0pt
\usepackage{bm}
\usepackage{bbm}
\usepackage{dsfont}
\usepackage{xspace}
\usepackage{cancel}
\usepackage{mathtools}
\usepackage{amsthm}
\makeatletter
\@ifundefined{theorem}{\newtheorem{theorem}{Theorem}}{}
\@ifundefined{proposition}{\newtheorem{proposition}[theorem]{Proposition}}{}
\makeatother
\DeclareMathOperator{\Cone}{Cone}
\DeclareMathOperator{\Len}{Len}
\newcommand{\QQ}{\mathbb{Q}}      % for Rationals
\newcommand{\RR}{\mathbb{R}}      % for Reals
\newcommand{\cF}{\mathcal{F}}    
\newcommand{\la}{\langle}    
\newcommand{\ra}{\rangle}   
\title[Bounds on Successive Minima of Orders in Number Fields and S]{Bounds on Successive Minima of Orders in Number Fields and Scrollar Invariants of Curves}
\author{Sameera Vemulapalli}
\date{Source: arXiv:2207.10522v2 (CC BY 4.0)}
\begin{document}
\maketitle

\noindent\emph{Source and licence.} Bounds on Successive Minima of Orders in Number Fields and Scrollar Invariants of Curves. Version arXiv:2207.10522v2 (CC BY 4.0), distributed under the Creative Commons Attribution 4.0 International licence, as recorded in the arXiv licence metadata for this exact version and shown on the abstract page https://arxiv.org/abs/2207.10522v2. Full rights notice: licenses/arxiv-2207.10522-CC-BY-4.0.txt.\par\smallskip

\noindent\textit{Editorial scope.} Counted unit: the Proposition whose source label is \texttt{lem:4p} in the article (source0.tex lines 2104--2110, source label \texttt{lem:4p}), delivered together with the article's own complete original proof of it (source0.tex lines 2111--2310, one \texttt{proof} environment of 200 lines). No mathematics is written, repaired or re-derived: statement and proof are the article's own text line for line. Required context retained uncounted: none. Every definition, display and label of those lines is retained unchanged. Omitted from this package and not redistributed: 1--2103, 2311--2553. Nothing inside a retained excerpt is omitted or altered: every retained body is delivered exactly as the article's lines are written, so no omission receipt of any kind accompanies this package. Citations: the retained excerpts carry no citation command at all, so no bibliography entry of the article is transcribed and no reference is redistributed. Reference adaptation: none was needed and none was made; every reference inside the delivered unit resolves inside it, so no unresolved reference or citation is produced. Printed slips kept literal: no printed word, symbol, hypothesis or number of the counted statement or of its proof was altered. Layout and class: the article's own class is \texttt{amsart} with packages including amssymb, amsthm, inputenc, geometry, pst-node, tikz-cd, thmtools, thm-restate, hyperref, amsrefs, cleveref, amsmath, amsfonts, mathtools; the wrapper is the lane's pinned \texttt{amsart} editorial wrapper at 10pt on A4, which restates the theorem-like environments the retained text needs and adds the article's own macro definitions that the retained text needs (\texttt{\textbackslash Cone}, \texttt{\textbackslash Len}, \texttt{\textbackslash QQ}, \texttt{\textbackslash RR}, \texttt{\textbackslash cF}, \texttt{\textbackslash la}, \texttt{\textbackslash ra}), with extra wrapper packages (bm, bbm, dsfont, xspace, cancel, mathtools). The delivered numbering is the unit's own. Assets: no retained span contains a figure environment, a graphics-inclusion command, a PDF-page inclusion, a table or a TikZ picture; no image file is copied into this package, the package declares no asset dependency and no image file is redistributed with it. Licence: version arXiv:2207.10522v2 of this article is distributed under the Creative Commons Attribution 4.0 International licence, as recorded in the arXiv licence metadata for this exact version and shown on the abstract page https://arxiv.org/abs/2207.10522v2. A full rights notice is delivered at licenses/arxiv-2207.10522-CC-BY-4.0.txt. Characters: every retained excerpt is pure ASCII, so the delivered engine, XeLaTeX with its default Latin Modern text font, reports no missing character.
\begin{proposition}
\label{lem:4p}
Let $n = 4p$ for $p$ a prime not equal to $2$ or $3$. Then there exists a flag $\cF$ of a degree $n$ number field such that 
\[
	P_{T_{\cF}} \not \subseteq \Len_{(2,2,p)} \cup \Len_{(2,p,2)} \cup \Len_{(p,2,2)}.
\]
\end{proposition}

\begin{proof}
It suffices to show that there exists a flag $\cF$ and a point 
\[
	\mathbf{x} \in \Cone(P_{T_{\cF}}) \setminus \Len_{(2,2,p)} \cup \Len_{(2,p,2)} \cup \Len_{(p,2,2)}.
\]

\noindent \textbf{Part A: defining the flag $\cF$.}
Choose $e_1,e_2,e_3 \in \overline{\QQ}$ such that:
\begin{itemize}
\item $e_1$ has degree $p$;
\item the element $e_2$ has degree $2$;
\item  the element $e_3$ has degree $2$;
\item and the compositum $\QQ(e_1,e_2,e_3)$ has degree $4p$. 
\end{itemize}
 
Define a basis $\{ v_0,\dots,v_{4p-1}\}$ of $K$ as follows. Set 
\begin{align*}
v_0 &\coloneqq 1 \\
v_1 &\coloneqq e_1 \\
v_2 &\coloneqq e_2 \\
v_3 &\coloneqq e_2 e_1 \\
v_4 &\coloneqq e_1^2 \\
v_5 &\coloneqq e_2 e_1^2 \\
v_6 &\coloneqq e_1^3 \\
v_7 &\coloneqq e_2e_1^3.
\end{align*}
For $8 \leq i < n$ and $1 \leq i' < n$, write $i' = i'_1 + i'_2p + i_3'2p$ in mixed radix notation with respect to $(p,2,2)$. Inductively define $v_i$ as follows. Choose $i'$ minimal such that $e_1^{i'_1}e_2^{i'_2}e_3^{i'_3} \notin \{v_0,\dots,v_{i-1}\}$. Set $v_i = e_1^{i'_1}e_2^{i'_2}e_3^{i'_3}$. Observe that for $i \geq 2p$, we have $i = i'$. Let $\cF$ be the corresponding flag.

\noindent \textbf{Part B: explicit description of $\cF$.}
Note that:
\begin{align*}
F_1 &= \QQ\la  1, e_1 \ra \\
F_3 &= \QQ( e_2) \QQ \la 1,e_1 \ra\\
F_5 &= \QQ( e_2)\QQ\la  1, e_1, e_1^2 \ra \\
F_7 &= \QQ( e_2)\QQ \la 1, e_1, e_1^2, e_1^3 \ra \\
F_{3p-1} &= \QQ( e_1 )  \QQ \la 1, e_2, e_3 \ra.
\end{align*}
Therefore
\[
	F_1F_{3p-1}  = F_{3p-1}
\]
\[
F_3F_3 = F_5
\]
\[
F_3F_5 = F_7.
\]

\noindent \textbf{Part C: defining $\mathbf{x}$ when $p = 5$.}
Suppose $p=5$. Then let $\mathbf{x} \in \RR^{19}$ be as follows:
\begin{align*}
x_1 &\coloneqq 1 \\
x_2,x &\coloneqq 1.4 \\
x_4,x_5 &\coloneqq 2 \\
x_6,x_7 &\coloneqq 3 \\
x_8,\dots,x_{14} &\coloneqq 4 \\
x_{15},\dots,x_{19} &\coloneqq 5.1.
\end{align*}

\noindent \textbf{Part D: showing that $\mathbf{x} \in \Cone(P_{T_{\cF}})$ and $x_8 > x_5 + x_3$ and $x_{15} > x_1 + x_{14}$ when $p = 5$.}
It is clear that $0 \leq x_1 \leq \dots \leq x_{n-1}$ and $x_8 > x_5 + x_3$ and $x_{15} > x_1 + x_{14}$. We now show that for all $1\leq i,j < n$, we have $x_{T_{\cF}(i,j)} \leq x_i + x_j$. Fix $i,j$ and let $k = T_{\cF}(i,j)$.

\noindent \textbf{Case 1: $i = 1$.} If $j = 1$ then $k = 4$ and 
\[
	x_4 = 2 \leq 2x_1.
\]
Observe that if $j = 2,3$ then $k = 4,5$, and
\[
	x_k = 2 \leq 1.4 + 1 \leq x_1 + x_j.
\]
If $j = 4,5$, then $k = 6,7$, and 
\[
	x_k = 3 \leq 2 + 1 = x_1 + x_j.
\]
If $j = 6,7$, then $k \leq 10$, and
\[
	x_k \leq 4 \leq 3 + 1 = x_1 + x_j.
\]
If $8 \leq j < 15$, then as $v_1 = e_1$ and $F_{14} = \la e_1 \ra \{ 1, e_2, e_3 \}$, we have that $j < k < 14$. Thus
\[
	x_k \leq 4 \leq 4 + 1 = x_1 + x_j.
\] 
If $15\leq j < n$ then
\[
	x_k \leq 5.1 \leq 5.1 + 1 \leq x_1 + x_j.
\]

\noindent \textbf{Case 2: $i = 2,3$.}
If $j = 2,3$, then $k = 4,5$ so 
\[
	x_k = 2 \leq 1.4 + 1.4 = x_i + x_j.
\]
If $j = 4,5$ then $k = 6,7$ so
\[
	x_k = 3 \leq 2 + 1.4 = x_i + x_j.
\]
If $j = 6,7$ then $k < 15$ so
\[
	x_k \leq 4 \leq 3 + 1.4 \leq x_i + x_j.
\]
If $8 \leq j < n$ then
\[
	x_k \leq 5.1 \leq 4 + 1.4 \leq x_i + x_j.
\]

\noindent \textbf{Case 3: $i = 4,5$.}
If $j < 10$ then because $F_9 = \QQ( e_1,e_2 )$ we have $k < 10$. Thus
\[
	x_k \leq 4 \leq 2 + 2 \leq x_i + x_j. 
\]
If $10 \leq j < n$ then
\[
	x_k \leq 5.1 \leq 4 + 2 \leq x_i + x_j.
\]

\noindent \textbf{Case 4: $i \geq 6$.}
Then 
\[
	x_k \leq 5.1 \leq 3 + 3 \leq x_i + x_j.
\]

\noindent \textbf{Part E: finishing the proof when $p = 5$.}
Because 
\[
	\Len_{(2,2,5)} \cup \Len_{(2,5,2)} \subset \{\mathbf{x} \in \RR^{19} : x_{15} \leq x_1 + x_{14} \}
\]
and
\[
	\Len_{(5,2,2)} \subset \{\mathbf{x} \in \RR^{19} : x_8 \leq x_5 + x_3 \},
\]
we have that 
\[
	\mathbf{x} \notin \Cone(\Len_{(2,2,5)} \cup \Len_{(2,5,2)} \cup \Len_{(5,2,2)})
\]
Hence, the proof is complete for $p = 5$.

\noindent \textbf{Part F: defining $\mathbf{x}$ when $p \neq 5$.}
If $p \neq 5$, let $\mathbf{x} \in \RR^{4p-1}$ be as follows.
\begin{align*}
x_1 &\coloneqq 1 \\
x_2, x_3 &\coloneqq 1.4 \\
x_4, x_5 &\coloneqq 2 \\
x_6, \dots, x_{3p-1} &\coloneqq 2.9 \\
x_{3p}, \dots, x_{4p-1} &\coloneqq 4.
\end{align*}

\noindent \textbf{Part G: showing that $\mathbf{x} \in \Cone(P_{T_{\cF}})$ and $x_6 > x_3 + x_3$ and $x_{3p} > x_1 + x_{3p-1}$.}
It is clear that $0\leq x_1 \leq \dots \leq x_{n-1}$ and $x_6 > x_3 + x_3$ and $x_{3p} > x_1 + x_{3p-1}$. We now show that for all integers $1 \leq i,j < n$, we have $x_{T_{\cF}(i,j)} \leq x_i + x_j$. Fix $i,j$ and let $k = T_{\cF}(i,j)$.

\noindent \textbf{Case 1: $i = 1$.}
If $j = 1$, then $k = 4$ and
\[
	x_4 = 2 \leq 2x_1.
\]
Observe that if $j = 2,3$, then $k = 4,5$, and
\[
	x_k \leq x_5 = 2 \leq 1.4 + 1 \leq x_1 + x_j.
\]
If $j = 4,5$, then $k = 6,7$, and
\[
	x_k \leq x_5 = 2.9 \leq 2 + 1 \leq x_1 + x_j.
\]
If $6 \leq j < 3p$, then as $v_1 = e_1$ and $F_{3p-1} = \QQ( e_1 ) \QQ \la 1, e_2, e_3 \ra$, we have $j < k < 3p$. Thus
\[
	x_k = 2.9 \leq 2.9 + 1 = x_j + x_1.
\]
If $j \geq 3p$, then
\[
	x_k = 4 \leq 4 + 1 \leq x_1 + x_j.
\]
\textbf{Case 2: $i = 2,3$.} 
If $j = 2,3$, then $k \leq 5$ and
\[
	x_k = 2 \leq 1.4 + 1.4 = x_i + x_j.
\]
If $j = 4,5$ then $k \leq 7$ and
\[
	x_k \leq 2.9 \leq 2 + 1.4 = x_i + x_j.
\]
If $j \geq 6$, then
\[
	x_k \leq 4 \leq 2.9 + 1.4 = x_i + x_j.
\]
\textbf{Case 3: $i \geq 4$.}
Then we have
\[
	x_k \leq 4 \leq 2 + 2 \leq x_i + x_j.
\]

\noindent \textbf{Part H: finishing the  proof when $p \neq 5$.}
By definition, we have
\[
	\Len_{(2,2,p)} \cup \Len_{(2,p,2)} \subset \{\mathbf{x} \in \RR^{4p-1} : x_{3p} \leq x_1 + x_{3p-1} \}
\]
and because $p \geq 7$, we have
\[
	\Len_{(p,2,2)} \subset \{\mathbf{x} \in \RR^{4p-1} : x_6 \leq x_3 + x_3 \},
\]
Therefore, $\mathbf{x} \notin \Cone(\Len_{(2,2,p)} \cup \Len_{(2,p,2)} \cup \Len_{(p,2,2)})$.
\end{proof}

\end{document}
